%% Question bank chapter file - generated by bank_tools.py build (do not edit while the bank is still being built:
%% it is regenerated from the per-test data in Question Bank\data). Every question is one bankq block:
%%   \begin{bankq}[id=..., source={...}, paper=.., qno=.., type=..., subtopic={...}, answer={...}, answer_src=key|solved,
%%                 working={...}, difficulty=1-10, diff_parts={K S I C F}, diff_note={...},
%%                 group=<passage id or empty>, also_in={duplicate copies in other tests}]
%% difficulty = min(10, K+S+I+C+F-1): K concepts 1-3, S steps 1-4, I insight 0-2, C computation 0-2, F format 0-1
%%   ...question LaTeX, options, figures (\bfig / \bfigi -> img/)...
%%   \end{bankq}
%% Passages shared by several questions are bankpassage blocks; their questions carry group=<passage id>.
\documentclass[10pt]{article}
\usepackage[a4paper,margin=1.8cm]{geometry}
\usepackage[T1]{fontenc}
\usepackage{amsmath,amssymb,mathtools}
\usepackage{graphicx}
\usepackage[export]{adjustbox}
\usepackage{xcolor,needspace,keyval}
\usepackage{enumitem}
\usepackage{booktabs,array,tabularx}
\usepackage[hidelinks]{hyperref}
\setlength{\parindent}{0pt}
\setlength{\parskip}{3pt}
\makeatletter
\define@key{bq}{id}{\def\bq@id{#1}}\define@key{bq}{source}{\def\bq@src{#1}}\define@key{bq}{paper}{\def\bq@paper{#1}}
\define@key{bq}{qno}{\def\bq@qno{#1}}\define@key{bq}{type}{\def\bq@type{#1}}\define@key{bq}{subtopic}{\def\bq@sub{#1}}
\define@key{bq}{answer}{\def\bq@ans{#1}}\define@key{bq}{answer_src}{\def\bq@asrc{#1}}\define@key{bq}{working}{\def\bq@work{#1}}
\define@key{bq}{group}{\def\bq@group{#1}}\define@key{bq}{also_in}{\def\bq@also{#1}}
\define@key{bq}{difficulty}{\def\bq@diff{#1}}\define@key{bq}{diff_parts}{\def\bq@dparts{#1}}\define@key{bq}{diff_note}{\def\bq@dnote{#1}}
\newenvironment{bankq}[1][]{\def\bq@work{}\def\bq@also{}\def\bq@group{}\def\bq@diff{?}\def\bq@dparts{}\def\bq@dnote{}\setkeys{bq}{#1}%
  \par\needspace{4\baselineskip}\medskip\noindent{\small\textbf{\bq@id}\hfill\color{gray}[\bq@src, Paper \bq@paper, Q\bq@qno]
  \bq@type\ $\cdot$ \bq@sub\ $\cdot$ \textbf{D\bq@diff}}\par\nopagebreak\relax}
  {\par{\small\color{teal}Answer: \textbf{\bq@ans}\ifx\bq@work\@empty\else\ (solved: \bq@work)\fi
  \ifx\bq@also\@empty\else\par\color{gray}Also in: \bq@also\fi}\par}
\define@key{bp}{id}{\def\bp@id{#1}}\define@key{bp}{source}{\def\bp@src{#1}}\define@key{bp}{questions}{\def\bp@qs{#1}}
\newenvironment{bankpassage}[1][]{\setkeys{bp}{#1}\par\needspace{4\baselineskip}\bigskip\noindent
  {\small\textbf{Passage \bp@id}\hfill\color{gray}[\bp@src; questions \bp@qs]}\par\nopagebreak\relax\itshape}{\par}
\makeatother
\newcommand{\bfig}[2]{\par\smallskip{\centering\includegraphics[width=#1pt,max width=0.62\linewidth,max height=0.32\textheight]{img/#2}\par}\smallskip}
\newcommand{\bfigi}[2]{\includegraphics[width=#1pt,max width=\linewidth,max height=0.32\textheight,valign=c]{img/#2}}
\newcommand{\bmissing}[1]{\textit{[figure missing in the source conversion: \detokenize{#1}]}}
\newenvironment{qoptions}{\begin{itemize}[leftmargin=1.8em,itemsep=1pt,topsep=2pt,label={}]}{\end{itemize}}
\newcommand{\opt}[2]{\item[(#1)] #2}
\providecommand{\fwstart}{}\providecommand{\fwend}{}\providecommand{\qfullsplit}{}
\newcolumntype{L}[1]{>{\raggedright\arraybackslash}p{#1}}
\newcolumntype{C}[1]{>{\centering\arraybackslash}p{#1}}
\newcolumntype{R}[1]{>{\raggedleft\arraybackslash}p{#1}}
\providecommand{\sectionhdr}[1]{}\providecommand{\subsectionhdr}[1]{}


\begin{document}
{\Large\bfseries Carboxylic Acids --- Question Bank}\par\smallskip
{\small Chemistry chapter C21. Questions from all converted test papers, grouped by source series. Answers: printed key where it matches the paper, otherwise solved (answer\_src=solved, with working). Exact repeats across papers appear once, with also\_in listing the other copies.}\par\medskip\hrule

\section*{Score-I papers}
\begin{bankq}[id=SI12-P1-CHEM-S1-Q02, source={Score-I - Test 12 (09 Mar 2022)}, paper=1, qno=2, type=multi, subtopic={Derivatives}, answer={A,B,C,D}, answer_src=key, difficulty=7, diff_parts={K3 S4 I0 C0 F1}, diff_note={Four different named reactions/products each to be recognised (acid derivatives, Ca acetate heating)}]
Which reactions are correctly matched with their product:
\begin{qoptions}
\opt{A}{\bfigi{158.4}{SI12_p1_chem_s1_q02_A.png}}
\opt{B}{\bfigi{134.9}{SI12_p1_chem_s1_q02_B.png}}
\opt{C}{\bfigi{136.4}{SI12_p1_chem_s1_q02_C.png}}
\opt{D}{(CH$_3$COO)$_2$Ca \bfigi{55.4}{SI12_p1_chem_s1_q02_4.png}}
\end{qoptions}
\end{bankq}

\section*{TOAS papers}
\begin{bankq}[id=TOAS04-P1-CHEM-S2-Q10, source={TOAS - Test 04 (15 Mar 2022)}, paper=1, qno=10, type=multi, subtopic={Derivatives}, answer={A,B,C}, answer_src=key, difficulty=6, diff_parts={K2 S3 I1 C0 F1}, diff_note={Kolbe electrolysis of the diacid gives bicyclo[2.2.2]octane; symmetry elements}]
\bfigi{199.6}{TOAS04_p1_chem_s2_q10_1.png}
\begin{qoptions}
\opt{A}{Product contain $C_3$ axis of symmetry}
\opt{B}{Product is a bicyclo compound}
\opt{C}{Product contain plane of symmetry}
\opt{D}{Product can show stereoisomerism}
\end{qoptions}
\end{bankq}
\begin{bankq}[id=TOAS07-P1-CHEM-S1-Q01, source={TOAS - Test 07 (01 May 2022)}, paper=1, qno=1, type=single, subtopic={Derivatives}, answer={D}, answer_src=key, difficulty=4, diff_parts={K1 S3 I1 C0 F0}, diff_note={LiAlD4 reduces the lactone and ketone: deuterium labelling in product P}]
\bfigi{141.1}{TOAS07_p1_chem_s1_q01_1.png}
\begin{qoptions}
\opt{A}{\bfigi{33.0}{TOAS07_p1_chem_s1_q01_A.png}}
\opt{B}{\bfigi{30.1}{TOAS07_p1_chem_s1_q01_B.png}}
\opt{C}{\bfigi{34.2}{TOAS07_p1_chem_s1_q01_C.png}}
\opt{D}{\bfigi{33.9}{TOAS07_p1_chem_s1_q01_D.png}}
\end{qoptions}
\end{bankq}
\begin{bankq}[id=TOAS07-P2-CHEM-S1-Q01, source={TOAS - Test 07 (01 May 2022)}, paper=2, qno=1, type=multi, subtopic={Derivatives}, answer={A,B,C,D}, answer_src=key, difficulty=5, diff_parts={K1 S4 I0 C0 F1}, diff_note={Four reactions of acids and derivatives with products}]
Which reactions are correctly matched with their product :
\begin{qoptions}
\opt{A}{\bfigi{127.1}{TOAS07_p2_chem_s1_q01_A.png}}
\opt{B}{\bfigi{124.9}{TOAS07_p2_chem_s1_q01_2.png}}
\opt{C}{\bfigi{123.9}{TOAS07_p2_chem_s1_q01_C.png}}
\opt{D}{$\mathrm{(CH_3COO)_2Ca}\xrightarrow{\Delta}$\bfigi{50.2}{TOAS07_p2_chem_s1_q01_4.png}}
\end{qoptions}
\end{bankq}
\begin{bankpassage}[id=TOAS07-P2-CHEM-S3-G1-2, source={TOAS - Test 07 (01 May 2022)}, questions={1-2}]
Heating of the compounds generally loose $\mathrm{CO_2}$ and $\mathrm{H_2O}$. The process is called decarboxylation and dehydration.
\end{bankpassage}
\begin{bankq}[id=TOAS07-P2-CHEM-S3-Q01, source={TOAS - Test 07 (01 May 2022)}, paper=2, qno=1, type=numerical, subtopic={Derivatives}, answer={3.00}, answer_src=key, difficulty=3, diff_parts={K1 S2 I0 C0 F1}, diff_note={Which dicarboxylic acids give cyclic anhydrides}, group=TOAS07-P2-CHEM-S3-G1-2]
Which of the following di-carboxylic acid on heating gives cyclic anhydride?\\
(i) Oxalic acid (ii) Succinic acid (iii) Glutaric acid (iv) Maleic acid (v) Fumaric acid
\end{bankq}
\begin{bankq}[id=TOAS07-P2-CHEM-S3-Q02, source={TOAS - Test 07 (01 May 2022)}, paper=2, qno=2, type=numerical, subtopic={Derivatives}, answer={2.00}, answer_src=key, difficulty=3, diff_parts={K1 S2 I0 C0 F1}, diff_note={gamma and delta hydroxy acids give lactones}, group=TOAS07-P2-CHEM-S3-G1-2]
Which of the following compounds on heating gives cyclic ester?\\
(i) $\gamma$-hydroxy acid \quad (ii) Pimelic acid \quad (iii) Urea \quad (iv) Benzoic acid \quad (v) $\delta$-hydroxy acid
\end{bankq}
\begin{bankq}[id=TOAS08-P1-CHEM-S1-Q05, source={TOAS - Test 08 (15 May 2022)}, paper=1, qno=5, type=multi, subtopic={Derivatives}, answer={A,B,C}, answer_src=key, difficulty=6, diff_parts={K2 S4 I0 C0 F1}, diff_note={Nitrile and acid-chloride routes to aldehydes: which reactions give RCHO}]
Which of the following reaction(s) produce aldehyde as a product :
\begin{qoptions}
\opt{A}{\bfigi{101.4}{TOAS08_p1_chem_s1_q05_A.png}}
\opt{B}{\bfigi{134.7}{TOAS08_p1_chem_s1_q05_B.png}}
\opt{C}{$\mathrm{CH_3-CN}\ \xrightarrow[\text{(ii) }\mathrm{H_3O^+}]{\text{(i) }\mathrm{SnCl_2}/\mathrm{HCl}}$}
\opt{D}{$\mathrm{CH_3MgBr+CH_3-C\equiv N}\ \xrightarrow{\mathrm{H^+/H_2O}}$}
\end{qoptions}
\end{bankq}
\begin{bankq}[id=TOAS08-P2-CHEM-S3-Q07, source={TOAS - Test 08 (15 May 2022)}, paper=2, qno=7, type=numerical, subtopic={Derivatives}, answer={4.00}, answer_src=key, difficulty=6, diff_parts={K2 S3 I1 C0 F1}, diff_note={Moles of CO2 released by the polyfunctional acid}]
How many moles of $\mathrm{CO_2}$ are released in following reaction :
\bfig{144.0}{TOAS08_p2_chem_s3_q07_1.png}
\end{bankq}
\begin{bankq}[id=TOAS09-P2-CHEM-S1-Q07, source={TOAS - Test 09 (30 May 2022)}, paper=2, qno=7, type=single, subtopic={Derivatives}, answer={C}, answer_src=key, difficulty=1, diff_parts={K1 S1 I0 C0 F0}, diff_note={Derivatives: standard approach for this subtopic; 1 concept(s), 1 step(s)}]
Which of the following gives alcohol as product :
\begin{qoptions}
\opt{A}{\bfigi{91.3}{TOAS09_p2_chem_s1_q07_A.png}}
\opt{B}{CH--CN $\xrightarrow[\mathrm{HOH}]{\mathrm{SnCl_2,\ HCl}}$}
\opt{C}{\bfigi{101.5}{TOAS09_p2_chem_s1_q07_C.png}}
\opt{D}{\bfigi{102.7}{TOAS09_p2_chem_s1_q07_D.png}}
\end{qoptions}
\end{bankq}
\begin{bankq}[id=TOAS09-P2-CHEM-S1-Q08, source={TOAS - Test 09 (30 May 2022)}, paper=2, qno=8, type=single, subtopic={Derivatives}, answer={C}, answer_src=key, difficulty=1, diff_parts={K1 S1 I0 C0 F0}, diff_note={Derivatives: standard approach for this subtopic; 1 concept(s), 1 step(s)}]
Select reaction from the following in which only CO$_2$ gas is released :
\begin{qoptions}
\opt{A}{\bfigi{93.8}{TOAS09_p2_chem_s1_q08_A.png}}
\opt{B}{\bfigi{125.2}{TOAS09_p2_chem_s1_q08_B.png}}
\opt{C}{\bfigi{98.0}{TOAS09_p2_chem_s1_q08_C.png}}
\opt{D}{\bfigi{129.8}{TOAS09_p2_chem_s1_q08_4.png}}
\end{qoptions}
\end{bankq}
\begin{bankq}[id=TOAS09-P2-CHEM-S1-Q10, source={TOAS - Test 09 (30 May 2022)}, paper=2, qno=10, type=single, subtopic={Derivatives}, answer={B}, answer_src=key, difficulty=2, diff_parts={K1 S2 I0 C0 F0}, diff_note={Derivatives: standard approach for this subtopic; 1 concept(s), 2 step(s)}]
Following conversion can not be carried out by using sequence
\[\mathrm{CH_3{-}CH_2{-}\overset{O}{\overset{\|}{C}}{-}OH\ \longrightarrow\ CH_3{-}CH_2{-}\overset{O}{\overset{\|}{C}}{-}H}\]
\begin{qoptions}
\opt{A}{(i) LiAlH$_4$ , PBr$_3$ , DMSO /NaHCO$_3$ / $\Delta$}
\opt{B}{(i) SOCl$_2$ (ii) CH$_3$MgX (excess)}
\opt{C}{(i) Ca(OH)$_2$ (ii) (HCOO)$_2$Ca/ Dry distillation}
\opt{D}{(i) LiAlH$_4$ (ii) Cu/300$^\circ$C}
\end{qoptions}
\end{bankq}
\end{document}
